\chapter{The software}\label{app:software}
\setlength{\emergencystretch}{3em}%

Every number in this book comes from one Python package, \texttt{privacy\_\allowbreak{}risk},
in the repository that accompanies it, and from the scripts beside it. This appendix
is a reader's map of that code. It says what each module computes, names its main
functions, gives the chapters that use it and the test file that checks it, and then
explains how to reproduce the evidence. It ends with the gaps a review of the code
found, which the book reports and does not repair, since the book branch does not
change code.

The map is written for two readers. A release committee member who wants to know
where a figure in a decision record came from can follow a module name to a chapter
and back. A researcher who wants to change an assumption and rerun a result can follow
a chapter to the function that computes it and the test that pins it down.

\section{Layout and installation}

The repository has five parts that matter to a reader of this book: the package in
\texttt{src/\allowbreak{}privacy\_\allowbreak{}risk/\allowbreak{}}, the scripts in \texttt{scripts/\allowbreak{}}, the tests in
\texttt{tests/\allowbreak{}}, the configurations in \texttt{configs/\allowbreak{}} and the recorded runs in the
directories whose names begin with \texttt{artifacts-}. The manuscripts in
\texttt{paper/\allowbreak{}} read their tables from \texttt{paper/\allowbreak{}generated/\allowbreak{}}, which a script copies
from a recorded run, so a table in a paper cannot disagree with the run behind it once
that step has been taken.

The package needs Python 3.10 to 3.12 and installs with its development extras from
the repository root. It depends on NumPy, pandas, SciPy and Matplotlib, and on version
0.5.0 of the \texttt{wooldridge} package, which supplies the two public microdata
files of Chapter~\ref{ch:public}. % source: pyproject.toml
No confidential data and no network access are needed, and no module reads real
confidential data. The installation defines four commands: \texttt{privacy-risk-run}
runs the paper's experiments and the governance gate, \texttt{privacy-risk-neural}
runs the neural release pipeline, \texttt{privacy-risk-neural-demo} runs a small
synthetic example of the proposal-and-audit method, and \texttt{sentry-r} drives the
proposal model of Chapter~\ref{ch:proposals}.

\section{The modules}

The package has 27 modules besides its initialiser. % source: src/privacy_risk/ (28 .py files including __init__.py)
They fall into three groups, shown in Tables~\ref{app:tab-core}, \ref{app:tab-case}
and~\ref{app:tab-prop}: the core of the method, which implements the definitions and
propositions of Parts~I to~IV; the case studies and the run machinery, which produce
the evidence of Part~V; and the model-proposed attack line with its audit, which
Chapter~\ref{ch:proposals} reports.

\subsection*{The core of the method}

Table~\ref{app:tab-core} lists the eleven modules that implement the method itself.
The first four build the objects of Chapter~\ref{ch:scenario}. \texttt{synthetic}
generates the simulated register of 200{,}000 units, the sampling design and the
protection mechanisms, and \texttt{keys} computes the sample and population key
frequencies. \texttt{linkage} holds the posterior linkage rule and the Monte Carlo
attack of the declared scenario, with the adversary's error model as a small class
whose \texttt{scaled} and \texttt{misspecified} methods produce the stronger and the
mistaken adversaries of Chapter~\ref{ch:falsification}. \texttt{risk} computes the
prosecutor, journalist and marketer measures of Chapter~\ref{ch:measures} and fits the
negative binomial model for population cell sizes.

The next three implement the evidence of Part~II. \texttt{uncertainty} gives the exact
one-sided Clopper--Pearson bounds, the number of attacks needed to certify a tolerance
with no success, the selection-corrected bound on the worst target and the split of
variance between simulation and design. \texttt{sequential} gives the beta--binomial
e-process and the confidence sequence it induces, the posterior threshold at which a
paid adversary acts, and the transport bound used for the temporal test.
\texttt{governance} holds the threshold record, the rule \texttt{decide}, the implied
tolerance and the programme risk budget, and \texttt{frontier} finds the nondominated
settings and fits the value curve from which a harm weight implies a tolerance.

The last three serve Part~IV and Part~VI. \texttt{composition} computes the sharp
bounds that marginal risks imply and runs the pooled intersection attack of
Chapter~\ref{ch:programmes}; \texttt{dp} states what a privacy budget permits as a
bound on posterior movement; and \texttt{validation} holds the falsification tests of
Chapter~\ref{ch:falsification}.

\begin{table}[ht]
\centering
\caption{Modules that implement the method. Chapters are those that report results
from the module; the test file is in \texttt{tests/\allowbreak{}}.}
\label{app:tab-core}
\footnotesize
\setlength{\tabcolsep}{3pt}
\begin{tabular}{@{}>{\raggedright\arraybackslash}p{2.0cm}>{\raggedright\arraybackslash}p{3.6cm}>{\raggedright\arraybackslash}p{4.1cm}>{\raggedright\arraybackslash}p{1.8cm}>{\raggedright\arraybackslash}p{2.8cm}@{}}
\toprule
Module & What it computes & Main functions & Chapters & Test file \\
\midrule
\texttt{synthetic} & Simulated register, sampling design, mechanisms, small-cell
suppression & \texttt{make\_\allowbreak{}register}, \texttt{apply\_\allowbreak{}mechanism},
\texttt{suppress\_\allowbreak{}small\_\allowbreak{}cells} & \ref{ch:scenario}, \ref{ch:simulation} &
\texttt{test\_\allowbreak{}synthetic} \\
\texttt{keys} & Sample and population key frequencies & \texttt{key\_\allowbreak{}frequencies},
\texttt{sample\_\allowbreak{}population\_\allowbreak{}frequencies} & \ref{ch:scenario} &
\texttt{test\_\allowbreak{}keys} \\
\texttt{linkage} & Posterior linkage and the Monte Carlo attack; error models &
\texttt{posterior\_\allowbreak{}linkage}, \texttt{attack}, \texttt{per\_\allowbreak{}target\_\allowbreak{}risk},
\texttt{ErrorModel} & \ref{ch:scenario}, \ref{ch:simulation}, \ref{ch:falsification}
& \texttt{test\_\allowbreak{}linkage} \\
\texttt{risk} & Prosecutor, journalist and marketer measures; frequency model &
\texttt{prosecutor\_\allowbreak{}risk}, \texttt{fit\_\allowbreak{}frequency\_\allowbreak{}model},
\texttt{journalist\_\allowbreak{}from\_\allowbreak{}fit} & \ref{ch:measures}, \ref{ch:falsification} &
\texttt{test\_\allowbreak{}risk} \\
\texttt{uncertainty} & Exact bounds, certification cost, selection-corrected worst
target, nested variance & \texttt{binomial\_\allowbreak{}upper},
\texttt{certification\_\allowbreak{}replicates}, \texttt{selected\_\allowbreak{}maximum\_\allowbreak{}bound},
\texttt{nested\_\allowbreak{}variance} & \ref{ch:estimate}, \ref{ch:selection} &
\texttt{test\_\allowbreak{}uncertainty} \\
\texttt{sequential} & E-process, confidence sequence, adversary's acting threshold,
transport bound & \texttt{anytime\_\allowbreak{}upper}, \texttt{attack\_\allowbreak{}threshold},
\texttt{transport\_\allowbreak{}bound\_\allowbreak{}kl}, \texttt{monitoring\_\allowbreak{}study} & \ref{ch:monitoring},
\ref{ch:falsification} & \texttt{test\_\allowbreak{}sequential} \\
\texttt{governance} & Threshold record, decision rule, implied tolerance, risk budget
& \texttt{GovernanceThreshold}, \texttt{decide}, \texttt{implied\_\allowbreak{}tolerance},
\texttt{risk\_\allowbreak{}budget} & \ref{ch:thresholds}, \ref{ch:tolerance}, \ref{ch:protocol} &
\texttt{test\_\allowbreak{}governance} \\
\texttt{frontier} & Nondominated settings, fitted value curve &
\texttt{nondominated}, \texttt{choose\_\allowbreak{}max\_\allowbreak{}utility}, \texttt{fit\_\allowbreak{}value\_\allowbreak{}curve} &
\ref{ch:tolerance} & \texttt{test\_\allowbreak{}frontier} \\
\texttt{composition} & Marginal bounds and the pooled intersection attack &
\texttt{marginal\_\allowbreak{}bounds}, \texttt{independent\_\allowbreak{}union}, \texttt{pooled\_\allowbreak{}attack},
\texttt{pooled\_\allowbreak{}attack\_\allowbreak{}from\_\allowbreak{}sets} & \ref{ch:programmes} &
\texttt{test\_\allowbreak{}composition} \\
\texttt{dp} & Posterior bound a privacy budget permits; basic composition &
\texttt{posterior\_\allowbreak{}bound}, \texttt{basic\_\allowbreak{}composition} & \ref{ch:dp} &
\texttt{test\_\allowbreak{}dp} \\
\texttt{validation} & The falsification tests & \texttt{calibration\_\allowbreak{}table},
\texttt{stress\_\allowbreak{}ratios}, \texttt{composition\_\allowbreak{}test},
\texttt{full\_\allowbreak{}validation\_\allowbreak{}report} & \ref{ch:falsification} &
\texttt{test\_\allowbreak{}validation} \\
\bottomrule
\end{tabular}
\end{table}

\subsection*{Case studies and the run}

Table~\ref{app:tab-case} lists the six modules that produce the evidence of Part~V and
assemble a run. \texttt{credit} builds the synthetic credit register of 20{,}000 firms
and 60 lenders over six monthly vintages and measures its five harms: borrower
identification, the predictive value of an alleged lending relationship, attribute
disclosure of arrears status, reconstruction of a lender's position from published
totals, and differencing across vintages. \texttt{evaluation} attacks the same
register as a security evaluation would: seven protection policies against several
attacks, scored by precision, recall, worst subgroup and calibration under a wrong
error model, with the search for the floor a tolerance implies, per release and
pooled, and the share of loans a rotated pseudonym still allows to be relinked.
\texttt{realdata} runs the demographic path on the two public files.

\texttt{experiments} calls the others and returns each number the long form reports,
keeping four kinds of work apart: proof checks, simulation checks, falsification tests
and the gate. \texttt{pipeline} runs the experiments, applies the falsification tests
and the gate, and writes the evidence; it is the entry point of
\texttt{privacy-risk-run}. \texttt{reporting} draws the figures, writes the LaTeX table
fragments and the file of numbers, and writes the manifest of SHA-256 hashes.

\begin{table}[ht]
\centering
\caption{Modules that produce the case-study evidence and assemble a run.}
\label{app:tab-case}
\footnotesize
\setlength{\tabcolsep}{3pt}
\begin{tabular}{@{}>{\raggedright\arraybackslash}p{2.0cm}>{\raggedright\arraybackslash}p{3.6cm}>{\raggedright\arraybackslash}p{4.1cm}>{\raggedright\arraybackslash}p{1.8cm}>{\raggedright\arraybackslash}p{2.8cm}@{}}
\toprule
Module & What it computes & Main functions & Chapters & Test file \\
\midrule
\texttt{credit} & Synthetic credit register and its five harms &
\texttt{make\_\allowbreak{}credit\_\allowbreak{}register}, \texttt{borrower\_\allowbreak{}identification},
\texttt{relationship\_\allowbreak{}disclosure}, \texttt{reconstruction\_\allowbreak{}from\_\allowbreak{}totals},
\texttt{differencing\_\allowbreak{}risk} & \ref{ch:credit}, \ref{ch:protocol} &
\texttt{test\_\allowbreak{}credit} \\
\texttt{evaluation} & Seven policies against the credit attacks; floor search;
pooled attack; relinking & \texttt{evaluate}, \texttt{identification\_\allowbreak{}attack},
\texttt{programme\_\allowbreak{}attack}, \texttt{relink\_\allowbreak{}share} & \ref{ch:credit},
\ref{ch:programmes}, \ref{ch:protocol} & \texttt{test\_\allowbreak{}evaluation} \\
\texttt{realdata} & Census 2000 extract and mortgage file as releases; pooled attack
on real keys & \texttt{load\_\allowbreak{}census2000}, \texttt{build\_\allowbreak{}release}, \texttt{attack},
\texttt{composition} & \ref{ch:public} & \texttt{test\_\allowbreak{}realdata} \\
\texttt{experiments} & Every experiment of the long form & \texttt{run\_\allowbreak{}all},
\texttt{programme}, \texttt{programme\_\allowbreak{}scaling}, \texttt{credit\_\allowbreak{}register} & all of
Part~V & \texttt{test\_\allowbreak{}experiments} \\
\texttt{pipeline} & Falsification, gate and evidence for one run & \texttt{validate},
\texttt{gate}, \texttt{run} & \ref{ch:falsification}, \ref{ch:protocol} &
\texttt{test\_\allowbreak{}pipeline} \\
\texttt{reporting} & Figures, table fragments, numbers, SHA-256 manifest &
\texttt{figures}, \texttt{paper\_\allowbreak{}tables}, \texttt{paper\_\allowbreak{}numbers},
\texttt{write\_\allowbreak{}evidence} & this appendix & \texttt{test\_\allowbreak{}reporting} \\
\bottomrule
\end{tabular}
\end{table}

\subsection*{Model-proposed attacks and the registered audit}

Table~\ref{app:tab-prop} lists the ten modules of the line that Chapter~\ref{ch:proposals}
reports and Chapter~\ref{ch:open} assesses. \texttt{adaptive\_\allowbreak{}audit} is the registered
sequential gate of the companion journal article: coordinates with their own error
shares, release chains that audit floors from the strictest down, the zero-success
budget, and the controls on the audit process that Chapter~\ref{ch:protocol} describes,
with targets derived from a committed secret and a beacon and bets recorded in a
hash-chained ledger before each outcome. \texttt{neural\_\allowbreak{}pipeline} is the package's main
entry point for the proposal line: a network fitted to development attacks proposes the
bets and starting floor, a plan is frozen with a digest, and the release chain decides
on fresh targets. \texttt{neural\_\allowbreak{}trust} is an earlier proposal-and-audit procedure that
supplies no input to the book's tables. \texttt{sentry\_\allowbreak{}r} holds the scenario schema,
the intake of language-model proposals, the review and evidence join, the deep and
logistic models, the shortlist freeze and the final audit. \texttt{provenance} is the
keyed, hash-chained ledger that records every stage of that line in order and debits
each audit against a programme error budget. \texttt{multiscale} adds a feature
dictionary, exact sensitivities of the network and an audit bound under a declared
distribution shift.

The remaining four are experiments around the line. \texttt{injector\_\allowbreak{}search} runs the
budget-matched search on public Census states and confirms the selected settings on
fresh draws. \texttt{llm\_\allowbreak{}pilot} collects proposals from a hosted model under a fixed
cost cap and never invents a response; \texttt{llm\_\allowbreak{}pilot\_\allowbreak{}local} does the same through
a local open-weight model. \texttt{redteam} turns penetration-test findings into
declared adversary capabilities and audits the release against the stronger view they
grant.

\begin{table}[ht]
\centering
\caption{Modules of the model-proposed attack line and the registered audit.}
\label{app:tab-prop}
\footnotesize
\setlength{\tabcolsep}{3pt}
\begin{tabular}{@{}>{\raggedright\arraybackslash}p{2.2cm}>{\raggedright\arraybackslash}p{3.5cm}>{\raggedright\arraybackslash}p{3.9cm}>{\raggedright\arraybackslash}p{1.6cm}>{\raggedright\arraybackslash}p{2.9cm}@{}}
\toprule
Module & What it computes & Main functions & Chapters & Test file \\
\midrule
\texttt{adaptive\_\allowbreak{}audit} & Registered sequential gate; release chains;
committed targets; ledger & \texttt{AdaptiveAudit}, \texttt{register\_\allowbreak{}release\_\allowbreak{}chain},
\texttt{zero\_\allowbreak{}success\_\allowbreak{}budget}, \texttt{target\_\allowbreak{}seed}, \texttt{AuditLedger} &
\ref{ch:monitoring}, \ref{ch:protocol}, \ref{ch:proposals} &
\texttt{test\_\allowbreak{}adaptive\_\allowbreak{}audit} \\
\texttt{neural\_\allowbreak{}pipeline} & Network proposes bets and start floor; frozen
plan; release chain & \texttt{run}, \texttt{fit}, \texttt{release\_\allowbreak{}chain} &
\ref{ch:proposals} & \texttt{test\_\allowbreak{}neural\_\allowbreak{}pipeline} \\
\texttt{neural\_\allowbreak{}trust} & Earlier proposal and simultaneous audit &
\texttt{certify}, \texttt{propose\_\allowbreak{}and\_\allowbreak{}audit}, \texttt{zero\_\allowbreak{}success\_\allowbreak{}budget} &
\ref{ch:proposals} & \texttt{test\_\allowbreak{}neural\_\allowbreak{}trust} \\
\texttt{sentry\_\allowbreak{}r} & Schema, language-model intake, review, models, freeze, final
audit & \texttt{stage\_\allowbreak{}llm}, \texttt{review\_\allowbreak{}proposals}, \texttt{DeepRiskNet},
\texttt{shortlist}, \texttt{final\_\allowbreak{}audit} & \ref{ch:proposals} &
\texttt{test\_\allowbreak{}sentry\_\allowbreak{}r}, \texttt{test\_\allowbreak{}injector\_\allowbreak{}integrity} \\
\texttt{provenance} & Keyed hash-chained ledger of every stage; error budget &
\texttt{ProvenanceLedger}, \texttt{record\_\allowbreak{}audit}, \texttt{verify} &
\ref{ch:proposals}, \ref{ch:protocol} & \texttt{test\_\allowbreak{}provenance},
\texttt{test\_\allowbreak{}injector\_\allowbreak{}integrity} \\
\texttt{multiscale} & Feature dictionary, network sensitivities, shifted audit &
\texttt{MultiscaleFeatures}, \texttt{shifted\_\allowbreak{}audit},
\texttt{zero\_\allowbreak{}success\_\allowbreak{}plan} & \ref{ch:proposals} & \texttt{test\_\allowbreak{}multiscale} \\
\texttt{injector\_\allowbreak{}search} & Budget-matched search on public states;
fresh confirmation & \texttt{select\_\allowbreak{}and\_\allowbreak{}measure}, \texttt{confirm},
\texttt{paired\_\allowbreak{}bootstrap} & \ref{ch:proposals}, \ref{ch:open} &
\texttt{test\_\allowbreak{}injector\_\allowbreak{}search} \\
\texttt{llm\_\allowbreak{}pilot} & Hosted-model proposal collection with a cost cap &
\texttt{collect}, \texttt{request\_\allowbreak{}payload} & \ref{ch:open} &
\texttt{test\_\allowbreak{}llm\_\allowbreak{}pilot} \\
\texttt{llm\_\allowbreak{}pilot\_\allowbreak{}local} & Local open-weight proposal collection &
\texttt{collect}, \texttt{response\_\allowbreak{}schema} & \ref{ch:open} &
\texttt{test\_\allowbreak{}llm\_\allowbreak{}pilot\_\allowbreak{}local} \\
\texttt{redteam} & Penetration-test findings as adversary capabilities; audit &
\texttt{adversary\_\allowbreak{}view}, \texttt{PenTestPlan}, \texttt{run} & \ref{ch:protocol} &
\texttt{test\_\allowbreak{}redteam} \\
\bottomrule
\end{tabular}
\end{table}

Three test files check code outside the package: \texttt{test\_\allowbreak{}linkage\_\allowbreak{}injector}
checks the executable linkage experiment in \texttt{scripts/\allowbreak{}}, and
\texttt{test\_\allowbreak{}colab\_\allowbreak{}setup} checks the helper that clones and synchronises the
repository from a notebook. The initialiser itself only sets a version string.

\section{The scripts}

The package computes; the scripts run studies, write tables and check manuscripts.
Table~\ref{app:tab-scripts} groups the 33 scripts in \texttt{scripts/\allowbreak{}} by what they do,
leaving aside the superseded scripts in \texttt{scripts/\allowbreak{}archive}.
% source: scripts/ (33 .py files)
The first group writes the evidence the book reports, and each of its scripts names
its output file at the top of its source. The second group checks that the manuscripts
agree with that evidence. The third covers the public-file and model-proposed attack
studies, and the fourth is housekeeping.

\begin{table}[ht]
\centering
\caption{The scripts in \texttt{scripts/\allowbreak{}}, grouped by purpose, with the chapters that
report what they produce.}
\label{app:tab-scripts}
\footnotesize
\setlength{\tabcolsep}{3pt}
\begin{tabular}{@{}>{\raggedright\arraybackslash}p{4.6cm}>{\raggedright\arraybackslash}p{7.6cm}>{\raggedright\arraybackslash}p{2.2cm}@{}}
\toprule
Script & What it does & Chapters \\
\midrule
\multicolumn{3}{@{}l}{\emph{Evidence for the book}}\\
\texttt{rebuild\_\allowbreak{}evidence} & Rewrites figures, tables, numbers and manifest from a
stored run without re-simulating; \texttt{--publish} copies tables and figures to
the manuscripts & all \\
\texttt{programme\_\allowbreak{}robustness} & Floors by adversary knowledge, tolerance, register
draw and reading noise; writes \texttt{artifacts-robustness/\allowbreak{}robustness.\allowbreak{}json} &
\ref{ch:credit}, \ref{ch:open} \\
\texttt{robustness\_\allowbreak{}table} & Tables of floors from that file & \ref{ch:credit},
\ref{ch:open} \\
\texttt{audit\_\allowbreak{}process\_\allowbreak{}attacks} & Seed grinding, shadow audits and post-hoc bets,
with and without controls & \ref{ch:protocol} \\
\texttt{adaptive\_\allowbreak{}audit\_\allowbreak{}chain} & The release chain on the credit register &
\ref{ch:monitoring}, \ref{ch:proposals} \\
\texttt{adaptive\_\allowbreak{}audit\_\allowbreak{}comparison} & Operating characteristics of the gate on
Bernoulli streams & \ref{ch:monitoring} \\
\texttt{estimand\_\allowbreak{}check} & Checks that two identification tables share one estimand
& \ref{ch:credit} \\
\texttt{worked\_\allowbreak{}examples} & Recomputes every number in the worked examples &
\ref{ch:protocol} \\
\texttt{neural\_\allowbreak{}results} & Figure and table of the network's predictions &
\ref{ch:proposals} \\
\addlinespace
\multicolumn{3}{@{}l}{\emph{Checks on the manuscripts}}\\
\texttt{check\_\allowbreak{}paper\_\allowbreak{}numbers} & Every run value appears in the long form; every
quoted decimal in the article has a source & all \\
\texttt{check\_\allowbreak{}float\_\allowbreak{}refs} & Every table, figure and algorithm is labelled and cited
& -- \\
\texttt{check\_\allowbreak{}realdata\_\allowbreak{}paper} & Archived public-data claims checked separately &
\ref{ch:public} \\
\addlinespace
\multicolumn{3}{@{}l}{\emph{Public files and model-proposed attacks}}\\
\texttt{realdata\_\allowbreak{}replication}, \texttt{realdata\_\allowbreak{}neural\_\allowbreak{}pilot},
\texttt{realdata\_\allowbreak{}neural\_\allowbreak{}comparison}, \texttt{plot\_\allowbreak{}realdata\_\allowbreak{}neural} &
Public-file replication with source hashes; state-held-out neural pilot and
comparison; their plot & \ref{ch:public}, \ref{ch:proposals} \\
\texttt{public\_\allowbreak{}injector\_\allowbreak{}search}, \texttt{report\_\allowbreak{}public\_\allowbreak{}search} & Budget-matched
search on 30 Census states and its tables & \ref{ch:proposals}, \ref{ch:open} \\
\texttt{collect\_\allowbreak{}llm\_\allowbreak{}pilot}, \texttt{collect\_\allowbreak{}llm\_\allowbreak{}pilot\_\allowbreak{}local},
\texttt{report\_\allowbreak{}llm\_\allowbreak{}local} & Language-model proposal collection, hosted and local, and
the local arm's table & \ref{ch:open} \\
\texttt{injector\_\allowbreak{}experiment}, \texttt{linkage\_\allowbreak{}injector\_\allowbreak{}experiment} & Synthetic
injector ablation; executable nearest-record linkage pilot & \ref{ch:proposals} \\
\texttt{sentry\_\allowbreak{}r\_\allowbreak{}example}, \texttt{sentry\_\allowbreak{}r\_\allowbreak{}benchmark} & Synthetic inputs and the
logistic against deep benchmark & \ref{ch:proposals} \\
\texttt{multiscale\_\allowbreak{}demo} & Training, sensitivity and shifted-audit demonstration &
\ref{ch:proposals} \\
\texttt{pentest\_\allowbreak{}demo} & Penetration-test findings on a public Census release &
\ref{ch:protocol} \\
\addlinespace
\multicolumn{3}{@{}l}{\emph{Housekeeping}}\\
\texttt{write\_\allowbreak{}manifest} & SHA-256 manifest of a list of files & this appendix \\
\texttt{build\_\allowbreak{}diagrams}, \texttt{build\_\allowbreak{}notebook}, \texttt{validate\_\allowbreak{}notebook} &
Editable diagrams; the Colab notebook and its check & -- \\
\texttt{make\_\allowbreak{}acmart}, \texttt{sync} & ACM version of the earlier manuscript;
repository synchronisation & -- \\
\bottomrule
\end{tabular}
\end{table}

\section{Reproducing the evidence}

One command regenerates everything the book and the paper print. From the repository
root, \texttt{python reproduce.py --full} runs the paper's configuration, rewrites the
tables, figures and number files, recomputes the numbers of the book's worked examples and
checks them against the chapter text, runs the test suite and builds the book;
\texttt{--fast} does the same with reduced replicates in a scratch directory and leaves the
committed files alone. The script names the committed files that a run changed. When the
paper's configuration was rerun on 2 November 2026 in the environment used to build this
book, the headline results agreed with the stored run to about $10^{-8}$ in the least
significant digits, and the figure files differed at the byte level. The numbers printed in
the text did not change. The run also added the uniqueness check of
Chapter~\ref{ch:falsification} and the four new governance fields to the stored results, which
the stored run predated. Byte-identical output across machines is therefore not something
the repository can promise; agreement of the printed numbers is.

The evidence for Parts~II to~V comes from one configuration, \texttt{configs/\allowbreak{}paper.\allowbreak{}json},
which fixes a seed, a register of 200{,}000 units with a 5 per cent sample, 45
mechanism settings, the credit register, and the bands of the falsification tests.
% source: configs/paper.json
A run with it takes a few minutes on one core. A smaller configuration,
\texttt{configs/\allowbreak{}base.\allowbreak{}json}, finishes in seconds and is what continuous integration
runs. The sequence below installs the package, reruns the paper's configuration,
publishes the tables, checks the manuscript against the run and runs the test suite.

\begin{lstlisting}[language={}]
python -m pip install -e ".[dev]"
privacy-risk-run --config configs/paper.json --output artifacts-paper
python scripts/rebuild_evidence.py --artifacts artifacts-paper \
    --publish paper/generated
python scripts/check_paper_numbers.py --artifacts artifacts-paper \
    --paper paper/bis_ifc_full.tex
python scripts/programme_robustness.py
python scripts/robustness_table.py
python scripts/audit_process_attacks.py
pytest
\end{lstlisting}

A run writes six things into its output directory: the configuration copied verbatim,
\texttt{results.\allowbreak{}json} with every experiment's output, \texttt{decision.\allowbreak{}json} with the
governance record, the figures, the LaTeX table fragments with a file of named numbers,
and \texttt{manifest.\allowbreak{}json}, which holds the SHA-256 hash of every other file. Figures
are written without creation dates and the results are strict JSON, so two runs of one
configuration produce identical files and an identical manifest.
% source: README.md "Reproducibility contract"; src/privacy_risk/reporting.py write_evidence
\texttt{rebuild\_\allowbreak{}evidence} rewrites the presentation from a stored run without
simulating again, which is the step to use when only a table's layout changes.

The manifest proves less than it may appear to. A hash shows that a file has changed
since the run that wrote the manifest. It carries no key and no trusted time stamp, so
anyone who can rewrite the files can rewrite the manifest as well; for evidence against
the operator the manifest's own hash has to be published or time-stamped outside the
operator's control. The script \texttt{write\_\allowbreak{}manifest} says this in its documentation.
% source: scripts/write_manifest.py docstring

The test suite has more tests than the 195 it held when this section was first
written. On 2 November 2026, run in the Linux environment used to build this book, 271
passed. Eight tests that need the \texttt{wooldridge} package failed until it was
installed. The earlier count, from 1 November, was 194 passed and one failed.
% source: pytest runs on the book branch, 1 and 2 November 2026
The one failure of 1 November, in the test that checks the shipped manifest against the
files it names, was attributed to Windows line endings in a working copy; the test passed in the
Linux run of 2 November. The suite's configuration also requires 85 per cent
line coverage of the package. % source: pyproject.toml addopts --cov-fail-under=85
Continuous integration runs the linter, the tests, the notebook, the fast pipeline,
the public-file scripts and both number checks on Python 3.10, 3.11 and 3.12, and a
second workflow reruns the paper's configuration on each push to the main branch that
touches code or configuration, publishes the tables and commits the refreshed
evidence. % source: .github/workflows/ci.yml, refresh-evidence.yml

The repository holds an older count of its own tests. Its validation note records 178
and later 183 passing tests for earlier snapshots, and says that a claimed 194 was not
reproduced at the time. % source: VALIDATION.md
The count of 194 passing tests reported here is from the current branch.

\section{Modules added by the mathematics and code extension}

The extension of November 2026 added modules that implement its corrections and its development plan, and the
revision of this book added the code that Chapters~18 to~22 print numbers from. Table~\ref{app:tab-ext} lists them.
They were written from the chapters' specifications and are tested against the checks the chapters print, not against
an earlier engine, so a figure the book takes from the extension's own run is reproduced in direction and size, and
usually to the printed digit, but not always; each chapter says where its numbers differ from the draft's.
\texttt{scripts/\allowbreak{}extension\_\allowbreak{}evidence.\allowbreak{}py} regenerates the pooled-attack table of Chapter~19, and
\texttt{scripts/\allowbreak{}chapter22\_\allowbreak{}evidence.\allowbreak{}py} regenerates the numbers of Chapter~22.

\begin{table}[ht]
\centering\small
\caption{Modules added or extended for Chapters 18 to 22. The test files are in \texttt{tests/}.}
\label{app:tab-ext}
\begin{tabular}{p{2.1cm}p{6.0cm}p{1.7cm}p{3.0cm}}
\toprule
Module & What it computes & Chapter & Test file\\
\midrule
\texttt{uncertainty} & Simultaneous bounds at $\alpha/M$; Beta posterior quantiles and beta-binomial mixtures, as planning figures & 22 & \texttt{test\_ext\_planner}\\
\texttt{evidence} & Measurement records as integer counts; refusal of the final-audit split & 18, 22 & \texttt{test\_ext\_learning}\\
\texttt{dynamics} & Changing and frozen panels, and the pooled attack using presence and absence & 19 & \texttt{test\_ext\_models}\\
\texttt{graph} & Generator checks, propagation, node risk and the maximum-likelihood generator & 19, 20 & \texttt{test\_ext\_models}\\
\texttt{transport} & The finite-volume solver under the exact shear; refuses unstable steps and never clips & 20 & \texttt{test\_ext\_transport}\\
\texttt{bayes} & Finite-prior posteriors, predictive mixture and the identifiability spectrum & 20 & \texttt{test\_ext\_transport}\\
\texttt{planner} & Count-likelihood planner with hand-written gradients, temperature calibration, grouped splits & 22 & \texttt{test\_ext\_planner}\\
\texttt{proposals} & Strict intake for proposed scenarios and plan hashes & 22 & \texttt{test\_ext\_planner}\\
\texttt{betting} & Predictable bets with a lower margin, log wealth and the sequential audit & 6, 18, 22 & \texttt{test\_ext\_learning}, \texttt{test\_ext\_planner}\\
\texttt{registry} & Programme budget, frozen plans, reservation before outcomes in SQLite & 18, 22 & \texttt{test\_ext\_learning}, \texttt{test\_ext\_planner}\\
\texttt{sensitivity} & Finite-time growth rates, the diagnostic horizon, decision flip rates and the stress gate & 21 & \texttt{test\_ext\_sensitivity}\\
\texttt{gate} & The release gate, which withholds rather than relaxes & 18 & \texttt{test\_ext\_learning}\\
\texttt{dp} & Attack ceiling, advanced composition, Laplace scale to $\varepsilon$, and the postprocessing inequality for a $[0,1]$-valued computation & 10, 21 & \texttt{test\_dp}, \texttt{test\_ext\_sensitivity}\\
\bottomrule
\end{tabular}
\end{table}

\paragraph{From equation to code.} Table~\ref{app:tab-eqmap} maps the main displayed results of Chapters~20 to~22 to the functions that compute them, so a reader can check one against the other.

\begin{table}[ht]
\centering\small
\caption{Results of Chapters 20 to 22 and the code that computes them.}
\label{app:tab-eqmap}
\begin{tabular}{p{4.6cm}p{5.2cm}p{3.0cm}}
\toprule
Result & Function & Chapter\\
\midrule
Upwind--diffusion update and stability bound $\Delta t\,R\le0.8$ & \texttt{transport.transport} & 20\\
Exact diffusion benchmark & \texttt{transport.exact\_diffusion} & 20\\
Finite-prior posterior and predictive mixture & \texttt{bayes.grid\_posterior}, \texttt{predictive\_pmf} & 20\\
Singular values of the sensitivity matrix & \texttt{bayes.sensitivity\_spectrum} & 20\\
Finite-time growth rate and diagnostic horizon & \texttt{sensitivity.finite\_time\_growth}, \texttt{diagnostic\_horizon} & 21\\
Flip rate and the stress gate & \texttt{sensitivity.decision\_flip\_rate}, \texttt{robust\_gate} & 21\\
Postprocessing inequality & \texttt{dp.postprocessing\_bound} & 21\\
Count loss and its gradient & \texttt{planner.objective} & 22\\
Temperature calibration & \texttt{planner.calibrate\_temperature} & 22\\
Credible against exact bound & \texttt{uncertainty.beta\_posterior\_upper}, \texttt{binomial\_upper} & 22\\
Predictable bets and the betting audit & \texttt{betting.predictable\_bets}, \texttt{run\_audit} & 22\\
Spending before outcomes & \texttt{registry.Registry.reserve} & 22\\
\bottomrule
\end{tabular}
\end{table}

\paragraph{Not in this repository.} The draft's command-line entries \texttt{drisk proposals} and \texttt{drisk registry}, and its hosted-model adapter, are not in this repository. The reviewer workflow, the dispatcher from approved proposals to evaluators, an external timestamp for the registry and an adapter for a real register are specified and not built.

\section{Gaps in the code}

A review of the repository before the book was written, and the work for the chapters
of Part~VI, found the following gaps. None affects a number the book reports, since
each number is checked against the stored run, but several affect what the code can
be trusted to do on its own, and an institution that adopts the code should close
them.

\paragraph{The governance record held six of the nine elements; it now holds all of them.}
\texttt{governance.GovernanceThreshold} first recorded tolerance, owner, expiry date,
adversary, harm class and error level. It now also carries the legal basis, the protected
parties, the access route and the programme horizon, which are the remaining three
elements stored as four fields. Called with \texttt{strict=True}, \texttt{decide} returns
INCOMPLETE, before any other status, when one of them is blank, because a number with no
named legal basis is not yet a decision. The default stays lenient so that earlier records
still load. \texttt{decide} still does not check that the bound it receives was computed
at the threshold's error level.

\paragraph{One tolerance for the programme.} Validation and the gate now judge the pooled
bound against one declared value, \texttt{programme.\allowbreak{}pooled\_\allowbreak{}tolerance}
(0.12 in the paper's configuration, defaulting to the release tolerance), and
\texttt{decision.\allowbreak{}json} records it once. The entry
\texttt{tau\_\allowbreak{}programme} (0.5) is the total of the union risk budget and is recorded
as \texttt{union\_\allowbreak{}risk\_\allowbreak{}budget}. An earlier version judged the same bound
against both values. % source: configs/paper.json; artifacts-paper/decision.json

\paragraph{Falsification bands.} The calibration band in the configuration is now 0.03,
down from 0.06, which a model with a misstated error scale met with a gap of 0.049. The
pipeline also requires the band to reject that misspecified model before the calibration
test counts as valid (check \texttt{calibration\_\allowbreak{}band\_\allowbreak{}has\_\allowbreak{}power}),
which test configurations can switch off with \texttt{require\_\allowbreak{}calibration\_\allowbreak{}power}. The
frequency test checks the mean population cell size, which the fitted law recovers to
0.055 cells. A second check now compares the modelled journalist uniqueness probability
with its truth, whose relative error reaches 25 per cent on the paper's register, against a
limit of 0.30 that was set after the errors were seen. % source: configs/paper.json; artifacts-paper/results.json validation; paper/generated/numbers.tex runjournalisterrorhigh
The composition and temporal tests are recorded beside the validity flag and in the
expiry check, as the long form intends.

\paragraph{Counts rebuilt from rates.} \texttt{uncertainty.\allowbreak{}selected\_\allowbreak{}maximum\_\allowbreak{}bound}
used to rebuild success counts by rounding rates times replicates. It now takes the
integer counts directly through its \texttt{successes} argument, and keeps the rounding only as a
fallback for callers that hold rates alone. A test checks that the two agree when the rates are exact counts and that
non-integer or out-of-range counts are refused.

\paragraph{An unstated normality assumption.} \texttt{uncertainty.\allowbreak{}nested\_\allowbreak{}variance}
gives a variance split from which a chi-square interval for the share of variance due to the design
is built. That interval assumes the outer estimates are roughly normal. The docstring now says so, and
adds that the assumption is poor when the risk is near zero or the outer replicates are few.

\paragraph{Version strings.} The package metadata and the package initialiser now both
say version 0.7.0. Archives of the repository named 0.7.1 are not matched by any tag.
A reader who cites a version should cite the commit.

\paragraph{Differential privacy.} \texttt{dp} covered the posterior bound and basic composition only. It now
also has \texttt{attack\_ceiling}, the ceiling $\min(1,e^{\varepsilon}\kappa+\delta)$ of Chapter~\ref{ch:dp},
\texttt{advanced\_composition}, which is never worse than the basic total, and \texttt{laplace\_epsilon}, which
reports $\varepsilon$ for a given noise scale and sensitivity. Noise of scale $1/\varepsilon$ on a count gives $\varepsilon$ when a
record is added or removed and $2\varepsilon$ when one is replaced, and the guarantee is per record: a borrower with several loans
needs a sensitivity that counts all of them. The code still does not implement a mechanism.

\paragraph{Line endings.} The one failing test, described above, and the long list of
files that show as modified in the Windows working copy, have the same cause and the
same remedy.

\paragraph{What is not built.} The registry service that a deployment needs, with
durable per-programme budgets, signing keys and external time stamps, is specified and
only partly built: \texttt{registry.Registry} keeps budgets and plan hashes in a local SQLite file, and has no signing keys or external time stamps. The adapter for a real register is described and not written. Dependencies
other than \texttt{wooldridge} are not pinned, so a fixed seed does not promise
bit-for-bit equality across library versions. % source: docs/DEPLOYMENT_AND_PENTEST_2026_10_23.md; docs/CLAIMS_EVIDENCE.md; VALIDATION.md
The pooled attack in \texttt{composition.\allowbreak{}pooled\_\allowbreak{}attack}, by contrast, was checked
against Proposition~\ref{prop:nopool}: a target absent from a release contributes zero
to that release's marginal risk and does not narrow the intersection, as the
proposition and the paper state.

\paragraph{Availability.} The papers describe the pipeline as open. At the time of
writing the repository is private. Either it is made public before the papers or this
book circulate, or the description changes.

\paragraph{Gaps and the claims they touch.} The gaps above and the fit of
Chapter~\ref{ch:tolerance} bear on specific claims. The governance record that omits three
elements limits the completeness claims of Chapters~\ref{ch:thresholds} and~\ref{ch:protocol};
the work to close it is to fill those fields from a signed form and test the strict check on
the record. The value-curve fit returns no standard errors, which limits the implied tolerance
of Chapter~\ref{ch:tolerance}; the work is to return the covariance and bootstrap over the
frontier points.

\section{Which gap touches which claim}\label{app:gapclaims}

Table~\ref{app:tab-gaps} sets each known gap against the claims it can affect and the work
that would close it. A reader who wants to rely on one claim need read only its row.

\begin{table}[t]
\centering
\caption{Gaps, the claims they touch, and the work that closes them.}\label{app:tab-gaps}
\small
\setlength{\tabcolsep}{3pt}
\begin{tabular}{@{}>{\raggedright\arraybackslash}p{3.0cm}>{\raggedright\arraybackslash}p{4.4cm}>{\raggedright\arraybackslash}p{4.6cm}@{}}
\toprule
Gap & Claims it affects & Work to close it \\
\midrule
Synthetic register with chosen parameters & Every floor (75, 150, none) in Chapters 1, 5 to 8,
13, 16, 18, 19 & Rerun the selection on the institution's own register \\
\addlinespace
Reading precision not measured & Floors in Chapter 13 and the insider result &
Appendix~\ref{app:b-precision}, then the sources the institution holds \\
\addlinespace
Frequency test had checked the mean cell size & The validity flag for the journalist measure &
Closed in code; limit of 0.30 set after the errors were seen, so first valid use is the next run \\
\addlinespace
Calibration band chosen after the wider band failed & Validity of the calibration test &
A run on fresh draws with the band fixed beforehand \\
\addlinespace
Governance record held six elements & Completeness of the decision record &
Closed in code, with \texttt{strict} mode; the book's six-versus-nine wording is the author's \\
\addlinespace
Registry service specified, not built & Claim that retries cannot silently spend the budget &
A service with authenticated writes and a stored reservation time \\
\addlinespace
Adapter for a real register not written & Every transfer claim in Chapter 23 &
An institution's adapter that emits only aggregates \\
\addlinespace
Composition gap measured only on synthetic and Census data & The ratio of 1.57 to 2.05 and 1.86 &
Measure the ratio on the institution's own programme \\
\bottomrule
\end{tabular}
\end{table}

